\chapter{Background and Related Work}

\section{Continuous Generative Flows}

A continuous generative model parameterises a time-dependent vector field
$v_\theta(x_t, t)$ such that integrating the corresponding ODE
\begin{equation}
  \frac{dx_t}{dt} = v_\theta(x_t, t)
\end{equation}
from $t=1$ (noise) to $t=0$ (data) produces samples from the learned
distribution. Neural ODEs~\cite{chen2018neural} established the general
framework in which a neural network defines this derivative and a numerical
solver computes the transformation. The number of times the network is called
during this integration is the NFE; minimising NFE while preserving quality
is the central efficiency objective.

\section{Flow Matching}

Flow Matching~\cite{lipman2023flow} trains a velocity network on a conditional
probability path without simulating the ODE during training. For data sample
$x_0$ drawn from the data distribution and noise $\varepsilon \sim \mathcal{N}(0,I)$,
the implementation defines a linear interpolation:
\begin{equation}
  x_t = (1 - t)\,x_0 + t\,\varepsilon, \qquad t \in [0, 1].
\end{equation}
The conditional target velocity is
\begin{equation}
  v = \varepsilon - x_0,
\end{equation}
and the training loss minimises the mean squared difference between the predicted
and target velocities:
\begin{equation}
  \mathcal{L}_{\mathrm{FM}}
    = \mathbb{E}_{x_0, \varepsilon, t}\bigl\|v_\theta(x_t, t)
      - (\varepsilon - x_0)\bigr\|_2^2.
\end{equation}
Sampling starts at $t=1$ and applies uniform backward Euler steps toward $t=0$:
\begin{equation}
  x_{t-h} = x_t - h\,v_\theta(x_t, t), \qquad h = 1/\mathrm{NFE}.
\end{equation}
One step can be too coarse when the learned field varies along the path;
more steps reduce integration error but increase latency.

\section{Logit-Normal Time Sampling}

The logit-normal variant draws the training time as
$t = \sigma(u)$ where $u \sim \mathcal{N}(0, 1)$ and $\sigma$ is the sigmoid
function, concentrating samples near $t = 0.5$~\cite{esser2024scaling}.
The path, target, backbone, and sampler remain unchanged; only the empirical
distribution of $t$ during training is altered. Perceptually important
intermediate noise levels receive more gradient signal, which can improve
sample quality without changing the model architecture.

\section{Mean Flow}

Mean Flow~\cite{geng2025mean} replaces instantaneous velocity with average
velocity over an interval. The network $u_\theta(z_t, r, t)$ models the
mean velocity over $[r, t]$, enabling the displacement rule:
\begin{equation}
  z_r = z_t - (t - r)\,u_\theta(z_t, r, t).
\end{equation}
Training uses the Mean Flow Identity, which relates interval-average velocity
to instantaneous velocity via a Jacobian-vector product (JVP):
\begin{equation}
  u(z_t, r, t)
    = v(z_t, t)
      - (t - r)\,\frac{d}{dt}\,u(z_t, r, t),
\end{equation}
where the total derivative $du/dt$ is evaluated along the direction
$(v, 0, 1)$ for $(z, r, t)$ using forward-mode automatic differentiation.
With probability 0.25 the endpoints coincide ($r = t$), recovering the
instantaneous velocity objective and regularising training.

At inference, setting $r = 0$ and $t = 1$ enables a direct noise-to-image
displacement in a single network call---the key one-step generation property.
Multiple steps are also possible by partitioning $[0,1]$ into sub-intervals.

The implementation used here is an educational CIFAR-10 adaptation;
it is not a reproduction of the paper's large-scale headline result.

\section{Rectified Flow and Reflow}

Rectified Flow~\cite{liu2022flow} learns straight trajectories between noise
and data, and its \emph{reflow} procedure distils a trained flow into a
single-step model by generating noise-data pairs using the original model and
training on those straight-line pairs. While not implemented in the submitted
system, this related approach contextualises MF's goal of one-step generation.

\section{Consistency Models}

Consistency Models~\cite{song2023consistency} train a network to map any point
on a trajectory directly to its endpoint, enforcing self-consistency across the
trajectory. Like MF, this enables single-step generation; unlike MF, it uses an
exponential moving average (EMA) teacher rather than a JVP-based identity.

\section{Backbone, Dataset, and Evaluation Metrics}

\paragraph{Backbone.}
U-Net-style encoders, decoders, and skip connections provide multi-scale
processing~\cite{ronneberger2015unet}. A sinusoidal time embedding is injected
into every residual block, making the architecture time-conditioned without
altering its structure between algorithms.

\paragraph{CIFAR-10.}
The dataset comprises 60{,}000 $32\!\times\!32$ RGB images in ten balanced
classes, conventionally split 50{,}000 train / 10{,}000 test~\cite{krizhevsky2009learning}.
Its manageable size makes repeated training, checkpoint sampling, and metric
computation practical on consumer hardware.

\paragraph{Fr\'echet Inception Distance.}
FID compares the multivariate Gaussian fitted to Inception-v3 features of real
and generated images~\cite{heusel2017gans}. Lower FID indicates closer
distributional alignment. Estimates from 1{,}000 samples have higher variance
than the standard 50{,}000-sample benchmark and should not be compared directly
to published results.

\paragraph{Inception Score.}
IS rewards both confident predictions (low entropy of $p(y|x)$) and diverse
predictions (high entropy of $p(y)$)~\cite{salimans2016improved}. It does not
compare against real images, making it an incomplete quality metric when used
alone.
